\section{Sistema de métricas: qué más observar además de la tasa de conexión entre personas reales}

Tristan menciona que la estrella polar es la tasa de conexión entre personas reales, y esa dirección es muy acertada. Sin embargo, un producto social no puede guiarse por una sola métrica. La tasa de conexión entre personas reales debe analizarse junto con la retención, la tasa de representación errónea, el grado de completitud del context, la tasa de revocación de privacidad, la proporción de interacciones IA/personas reales y la calidad de las relaciones a largo plazo.

\subsection{Tabla de métricas}

\noindent\begin{longtable}{p{0.24\textwidth}p{0.32\textwidth}p{0.35\textwidth}}
\toprule
Métrica & Qué mide & Riesgo de mal uso \\
\midrule
Tasa de conexión entre personas reales & Proporción de interacciones efectivas entre personas reales propiciadas por la IA & Puede inducir a molestar en exceso. \\
Proporción de actividad de personas reales & Porcentaje de las interacciones de la red que hacen personas reales y no la IA & Si es demasiado baja, la red se vuelve una IA que habla consigo misma. \\
Grado de completitud del context & Si el usuario proporciona suficiente context utilizable & Perseguirlo en exceso invade la privacidad. \\
Tasa de representación errónea & Proporción de salidas del doble digital que el usuario rechaza & Requiere un registro minucioso y una corrección rápida. \\
Retención & Si el usuario sigue regresando & Puede inflarse con una actividad aparente que no es real. \\
Calidad de las relaciones & Si las conexiones son duraderas, valiosas y de bajo daño & Es difícil de cuantificar, pero es la más cercana a la visión. \\
\bottomrule
\end{longtable}

\begin{knowledgebox}{Por qué las métricas deben analizarse en conjunto}
Fijarse solo en la tasa de conexión entre personas reales puede fomentar que se atraiga a la gente de forma brusca; fijarse solo en la retención puede fomentar la adicción; fijarse solo en el número de interacciones con la IA puede generar una prosperidad ficticia. La IA social necesita un conjunto de métricas que se equilibren entre sí.
\end{knowledgebox}

\subsection{Resumen de la sección}

La tasa de conexión entre personas reales es la estrella polar, pero no el único tablero de control. Un buen sistema de métricas para la IA social debe considerar a la vez la conexión, la privacidad, la representación errónea y la calidad de las relaciones.
